% ======================================================================
% Section 3 -- Preliminaries (~1 page budget)
% Say "neighborhoods", never "relations", for message-passing routes.
% ======================================================================
\section{Preliminaries}
\label{sec:preliminaries}

Three objects underlie the paper: the data structure (complexes and
cells), the models (GCCNs and neighborhoods), and the attribution
machinery (games and Shapley values). The molecule is the running
example: atoms are rank-0 cells, bonds rank-1, rings rank-2. A
glossary of every recurring term is in \cref{app:glossary}.

\subsection{Combinatorial Complexes}
\label{sec:prelim-cc}

\begin{definition}[Combinatorial complex \citep{hajij2023tdl}]
\label{def:cc}
A \emph{combinatorial complex} is a triple $(V, \mathcal{X}, \rk)$ where
$V$ is a finite set of vertices, $\mathcal{X} \subseteq
2^{V}\setminus\{\emptyset\}$ is a set of \emph{cells}, and $\rk :
\mathcal{X} \to \mathbb{Z}_{\ge 0}$ is a rank function monotone with
respect to inclusion: $x \subseteq y \implies \rk(x) \le \rk(y)$.
\end{definition}

Cells of rank $0$ and $1$ are nodes and edges; cells of rank $2$ are
faces, and higher ranks hold further multi-way structure. Graphs are
the rank-$\le 1$ special case. Hypergraphs carry set-type relations
and no rank-2 cells, so nothing in what follows assumes rank-2 cells
exist. A \emph{lifting} converts a graph into a complex by promoting
recognizable substructures to higher-rank cells, \eg chordless cycles
to faces \citep{bodnar2021cwn}; MANTRA
\citep{ballester2025mantra} is natively higher-order and needs no
lifting. We write $\mathcal{X}_r$ for the cells of rank $r$ and
$\mathbf{H}_r$ for the feature matrix on rank $r$.

\subsection{GCCNs and Neighborhoods}
\label{sec:prelim-gccn}

Message passing on a complex is organized by \emph{neighborhoods}:
routes along which cells exchange messages, each named by who talks to
whom \citep{papillon2025topotune}: \emph{incidences} $B_{r,r'}$
(\eg nodes to the edges containing them), \emph{up-adjacencies}
$A^{\uparrow}_{r}$ (\eg nodes joined by an edge, written
\nbhd{up\_adjacency-0}), and \emph{down-adjacencies}
$A^{\downarrow}_{r}$. A generalized combinatorial complex network
(GCCN) is specified by a set $\mathcal{N} = \{N_1, \dots, N_m\}$ of
neighborhoods; each layer runs one message-passing submodule per
neighborhood and aggregates:
\begin{equation}
\label{eq:gccn-layer}
\mathbf{h}^{(\ell+1)}_x
  \;=\;
  \phi\!\Bigl(\mathbf{h}^{(\ell)}_x,\;
  \bigotimes_{N \in \mathcal{N}}\;
  \bigoplus_{y \in N(x)}
  \psi_{N}\bigl(\mathbf{h}^{(\ell)}_x, \mathbf{h}^{(\ell)}_y\bigr)\Bigr),
\end{equation}
with $\bigoplus$ aggregating within and $\bigotimes$ across
neighborhoods. In TopoBench the candidate set for rank-$\le 2$
complexes has \unfrozen{\ResNumNeighborhoods} neighborhoods; which
subset to activate is ordinarily answered by sweeps that train one
model per candidate subset \citep{papillon2025topotune}.
\Cref{sec:performance} answers it with one training run of the
all-neighborhoods model plus attribution.

\subsection{Shapley Values and Interactions}
\label{sec:prelim-shapley}

A \emph{cooperative game} $(P, v)$ has players $P$, $|P| = n$, and a
value function $v : 2^{P} \to \mathbb{R}$, $v(\emptyset) = 0$, scoring
every \emph{coalition} (subset). The \emph{Shapley value}
\citep{shapley1953value} of player $i$ averages its marginal
contribution over all coalitions:
\begin{equation}
\label{eq:shapley}
\phi_i(v) \;=\; \sum_{S \subseteq P \setminus \{i\}}
  \frac{|S|!\,(n - |S| - 1)!}{n!}
  \bigl[v(S \cup \{i\}) - v(S)\bigr].
\end{equation}
It is the unique attribution satisfying \emph{efficiency}
($\sum_i \phi_i = v(P)$), \emph{symmetry}, \emph{null player}
($\phi_i = 0$ if $i$ never changes any coalition's value), and
\emph{linearity}. The \emph{interaction index}
\citep{grabisch1999axiomatic} extends \cref{eq:shapley} to coalitions
$T$ via discrete derivatives:
\begin{equation}
\label{eq:interaction}
I(T) \;=\; \sum_{S \subseteq P \setminus T}
  \frac{|S|!\,(n - |S| - |T|)!}{(n - |T| + 1)!}\,
  \Delta_{T} v(S),
\qquad
\Delta_{T} v(S) = \sum_{L \subseteq T} (-1)^{|T| - |L|}\, v(S \cup L),
\end{equation}
so $I(\{i,j\}) > 0$ indicates synergy and $I(\{i,j\}) < 0$ redundancy.
Exact computation needs $2^{n}$ evaluations; \cref{sec:framework}
gives the regimes where our games afford it and the sampled estimator
otherwise. Because the axioms are equalities a finite computation can
check, they double as executable tests of the model-plus-explainer
stack (\cref{app:experimental}).
